\section*{الفصل \ref{ch:g12:continental-crust} --- القشرة القارية}

\begin{solution}{exo:g12:continental-crust:1}
\omterm{def:g12:continental-crust:continental-crust}{القشرة القارية}: من 30 إلى \qty{70}{km}، ونحو \qty{41}{km} في المتوسط؛
صخور غرانيتية، ونيس وصخور متحولة \omterm{def:g11:igneous-rocks:volcanic-plutonic}{وصخور جوفية} أخرى تحت
غطاء رسوبي؛ ونحو \qty{61}{\percent} من السيليكا في المتوسط؛ ونحو
\qty{2830}{kg/m^3}؛ وحتى 4.03 مليار سنة (والحبيبات 4.4). \omterm{def:g11:ocean-ridges:oceanic-lithosphere}{القشرة المحيطية}:
نحو \qty{7}{km}؛ بازلت \omterm{def:g11:ocean-ridges:gabbro}{وغابرو}؛ ونحو \qty{50}{\percent} من السيليكا؛ ونحو
\qty{3000}{kg/m^3}؛ وأقل من 200 مليون سنة.
\end{solution}

\begin{solution}{exo:g12:continental-crust:2}
\qty{60.6}{\percent}: متوسطة، \omterm{def:g11:igneous-rocks:rock-names}{كالديوريت}. \qty{66.6}{\percent}:
\omterm{def:g11:igneous-rocks:felsic-mafic}{فلسية}، كالغرانيت.
\end{solution}

\begin{solution}{exo:g12:continental-crust:3}
\omterm{def:g11:ocean-ridges:oceanic-lithosphere}{القشرة المحيطية}، نحو \qty{3000}{kg/m^3}، تستقر على وشاح \omterm{def:g11:lithosphere-asthenosphere:lithosphere}{الغلاف الصخري}،
\omterm{def:g11:ocean-ridges:oceanic-lithosphere}{والغلاف الصخري المحيطي} القديم البارد كله كثيف بما يكفي ليغوص.
أما \omterm{def:g12:continental-crust:continental-crust}{القشرة القارية}، نحو \qty{2830}{kg/m^3}، فهي أخف بكثير من الوشاح
(نحو \qty{3380}{kg/m^3})، وطفوها شديد: فالقارة التي تصل إلى خندق
توقف \omterm{def:g11:subduction-zones:subduction}{الاندساس}.
\end{solution}

\begin{solution}{exo:g12:continental-crust:4}
(1) نواة قديمة مستقرة لقارة، لم تعد تتشوه. (2) أكثر
من مليار سنة. (3) صخور متحولة \omterm{def:g11:igneous-rocks:volcanic-plutonic}{وصخور جوفية}، هي \omterm{def:g12:continental-crust:root}{الجذور العميقة}
لجبال سابقة.
\end{solution}

\begin{solution}{exo:g12:continental-crust:5}
$4030/190 \approx 21$ مرة. فقاع المحيط يُصنع عند الظهرات
ويُدمَّر في نطاقات \omterm{def:g11:subduction-zones:subduction}{الاندساس}: ولا يبقى شيء منه أكثر من 200 مليون
سنة بكثير.
\end{solution}

\begin{solution}{exo:g12:continental-crust:6}
(1) $0.306 \times 210.4 \approx \num{64.4}\text{ مليون }\unit{km^2}$.
(2) $210.4 - 64.4 = \num{146.0}\text{ مليون }\unit{km^2}$.
(3) $4\pi \times 6371^2 \approx \qty{5.10e8}{km^2}$؛ $210.4/510.1 \approx
0.41$: أي \qty{41}{\percent}.
\end{solution}

\begin{solution}{exo:g12:continental-crust:7}
(1) $41100 \times 2830 \approx \qty{1.16e8}{kg}$. (2) $7100 \times 3000 =
\qty{2.13e7}{kg}$. (3) نحو 5.5 مرة.
\end{solution}

\begin{solution}{exo:g12:continental-crust:8}
\omterm{def:g12:continental-crust:root}{الجذر}: $2.5 \times 2830/550 \approx \qty{12.9}{km}$. والقشرة أسمك بمقدار
$2.5 \times 3380/550 \approx \qty{15.4}{km}$: أي نحو \qty{50}{km} إجمالًا،
وهذا ضمن المدى من 50 إلى \qty{70}{km} المقيس تحت الأحزمة الجبلية.
\end{solution}

\begin{solution}{exo:g12:continental-crust:9}
(1) نحو $0.103 + 0.100 + 0.102 \approx 0.30$: أي \qty{30}{\percent}.
(2) الشريط الأحدث، من 0 إلى 10 ملايين سنة (نحو \qty{10}{\percent})،
ويكاد يتعادل مع الشريط من 20 إلى 30. (3) لقد دُمّر القاع القديم في نطاقات
\omterm{def:g11:subduction-zones:subduction}{الاندساس}؛ وكلما كان القاع أقدم، كان ما زال منه أكثر.
\end{solution}

\begin{solution}{exo:g12:continental-crust:10}
$120 = 0.694\,h - 0.306 \times 100$، إذن $h = 150.6/0.694 \approx
\qty{217}{m}$. فالاقتصار على اليابسة يُسقط الحواف المنخفضة المغمورة
للقارات، فيخرج المتوسط أعلى بكثير.
\end{solution}

\begin{solution}{exo:g12:continental-crust:11}
الجزيرة: \omterm{def:g11:ocean-ridges:gabbro}{غابرو} تحت لابات \omterm{def:g11:seismic-structure:moho}{وموهو} على عمق \qty{8}{km}، فهي \omterm{def:g11:ocean-ridges:oceanic-lithosphere}{قشرة محيطية}
(وارتفاع الجزيرة لا يغير ذلك). والسهل: نيس \omterm{def:g11:seismic-structure:moho}{وموهو} على
عمق \qty{38}{km}، فهو \omterm{def:g12:continental-crust:continental-crust}{قشرة قارية}.
\end{solution}

\begin{solution}{exo:g12:continental-crust:12}
(1) الدقيق: $\tfrac{4}{3}\pi(6371^3 - 6329.9^3) \approx
\qty{2.0829e10}{km^3}$؛ والتقريبي: $4\pi \times 6371^2 \times 41.1
\approx \qty{2.0964e10}{km^3}$. (2) التقريب أكبر مما ينبغي بنحو
\qty{0.65}{\percent}، وهذا قريب من $H/R \approx 0.0065$: فالخطأ في طبقة
رقيقة مهمل.
\end{solution}

\begin{solution}{exo:g12:continental-crust:13}
(1) $60.6 = (66.6 + 2x)/3$، إذن $x = (181.8 - 66.6)/2 = \qty{57.6}{\percent}$.
(2) متوسطة. (3) \omterm{def:g12:continental-crust:continental-crust}{القشرة القارية} العميقة أفقر بالسيليكا من
القشرة العليا: فيها معادن \omterm{def:g11:igneous-rocks:felsic-mafic}{مافية} أكثر، وهي أغنى بالحديد والمغنيسيوم والكالسيوم.
\end{solution}

\begin{solution}{exo:g12:continental-crust:14}
(1) $3 \times 10^8/3 = 10^8$ سنة، أي 100 مليون سنة. (2) هي من رتبة
أعمار قاع المحيط نفسها: ثلاثة أرباعه أحدث من 100 مليون سنة،
ولا شيء منه أقدم من نحو 190. (3) يجب أن يُدمَّر بالسرعة التي يُصنع بها:
فهو يغوص عائدًا إلى الوشاح في نطاقات \omterm{def:g11:subduction-zones:subduction}{الاندساس}.
\end{solution}

\begin{solution}{exo:g12:continental-crust:15}
(1) $8.6 \times 10^9 / 4.0 \times 10^9 \approx \qty{2.2}{km^3}$ في السنة.
(2) المعدل الثابت ليس إلا احتمالًا واحدًا: فربما نمت القشرة سريعًا في البداية
ثم أُعيد تدويرها منذئذ، أو نمت متأخرة. وأقدم الصخور والحبيبات تثبت
أن بعض القشرة وُجد مبكرًا؛ لكنها لا تستطيع أن تخبر بمقداره، لأن القشرة القديمة
قد تتآكل أو تُحمل عائدة إلى الوشاح فلا تترك إلا آثارًا قليلة.
\end{solution}

\begin{solution}{pb:g12:continental-crust:1}
\textbf{1.} $4\pi \times 6371^2 \approx \qty{5.10e8}{km^2}$.

\textbf{2.} $210.4/510.1 \approx 0.41$: أي \qty{41}{\percent}.

\textbf{3.} $0.306 \times 210.4 \approx \num{64.4}\text{ مليون }\unit{km^2}$.

\textbf{4.} $510.1 - 210.4 \approx \num{299.7}\text{ مليون }\unit{km^2}$،
أي \qty{59}{\percent}.

\textbf{5.} تقع حواف القارات تحت بحار ضحلة، على
الأرصفة القارية.

\textbf{6.} السمك، \qty{41}{km}، صغير جدًا مقارنةً بنصف
القطر، \qty{6371}{km}: فالتقوس يغيّر الحجم بأقل من واحد
في المئة.

\textbf{7.} $210.4 \times 10^6 \times 41.1 \approx \qty{8.65e9}{km^3}$.

\textbf{8.} \qty{8.65e18}{m^3}.

\textbf{9.} $210.4 \times 10^6 \times 36 \approx \qty{7.57e9}{km^3}$.

\textbf{10.} $(8.65 - 7.57)/7.57 \approx 0.14$: أي \qty{14}{\percent}.

\textbf{11.} $8.65 \times 10^{18} \times 2830 \approx \qty{2.45e22}{kg}$.

\textbf{12.} $41100 \times 2830 \approx \qty{1.16e8}{kg}$.

\textbf{13.} $299.7 \times 10^6 \times 7.1 \approx \qty{2.13e9}{km^3}$، أي
\qty{2.13e18}{m^3}؛ وبالضرب في 3000: نحو \qty{6.4e21}{kg}.

\textbf{14.} \omterm{def:g12:continental-crust:continental-crust}{القشرة القارية}، أثقل بنحو 3.8 مرة: فهي أسمك بنحو ست
مرات، وهذا يفوق صغر مساحتها وكثافتها.

\textbf{15.} إنها مكوّنة من صخور غرانيتية، غنية بالكوارتز والفلسبارات، وهي
أخف المعادن المكوّنة للصخور؛ أما \omterm{def:g11:ocean-ridges:oceanic-lithosphere}{القشرة المحيطية} فمكوّنة من البازلت
\omterm{def:g11:ocean-ridges:gabbro}{والغابرو}، الغنيين بالبيروكسين والأوليفين.

\textbf{16.} $2.45 \times 10^{22}/5.97 \times 10^{24} \approx 0.0041$.

\textbf{17.} $3 \times 0.41 \times (41.1/6371) \times (2830/5513) \approx
0.0041$: النتيجة نفسها.

\textbf{18.} $6.4 \times 10^{21}/5.97 \times 10^{24} \approx 0.0011$: أي نحو
عُشر في المئة.

\textbf{19.} نحو \qty{0.4}{\percent}.

\textbf{20.} \omterm{def:g12:continental-crust:continental-crust}{القشرة القارية} نحو \textbf{\qty{0.4}{\percent} من كتلة
الأرض}.
\end{solution}
